\section{Rational packets and the split half-power}
\label{sec:rational-packets}

Let \(S_n=V(\sum x_i^3)\) over \(\Q\).  For the certified smooth rows
\(n=2,3,4\), define
\[
\mathsf E_n=\one\oplus(S_n,\pi_{\mathrm{prim}}^{S_n},n-1),
\qquad
\mathsf O_n=(X_{n,0},\pi_{2n-3}^{X_n},n-2),
\]
and
\[
\mathsf W_n=\mathsf E_n\oplus\mathsf O_n.
\]
The projectors are the ordinary hyperplane projectors and their middle
complements.  For the fivefold \(X_{4,0}\), for example,
\begin{equation}
\pi_{2i}=\frac16h^{5-i}\times h^i,\qquad
\pi_5=\Delta-\sum_{i=0}^5\pi_{2i}.
\label{eq:pi5}
\end{equation}
The coefficient \(1/6\) is the inverse degree of the complete
intersection.  No multiplicative Chow--K\"unneth assertion is used.

The source cohomology calculation gives
\begin{equation}
\rank\mathsf E_n=\frac{4^n+5}{3},\qquad
\rank\mathsf O_n=\frac{2(4^n-4)}3.
\label{eq:packet-ranks}
\end{equation}
Both summands are now defined over \(\Q\): the Fermat cubic was rational
from the outset, and Theorem~\ref{thm:family-descent} supplies
\(X_{n,0}\).  Base change preserves the projectors and ranks.  Adding
\eqref{eq:packet-ranks} proves
\[
\rank\mathsf W_n=4^n-1,
\]
which is \(15,63,255\) in the three certified rows.

\subsection{Placewise quadratic base change}

For a rational prime \(p\) outside the common bad set, write
\[
L_{K,p}(\mathsf W_{n,K},u)
=\prod_{v\mid p}L_{K,v}(\mathsf W_{n,K},u).
\]
If \(p\) splits in \(K\), both residue fields equal \(\F_p\).
Because \(\mathsf W_{n,K}\) is the base change of \(\mathsf W_n\), the
two local factors are identical:
\begin{equation}
L_{K,p}(\mathsf W_{n,K},u)
=L_{\Q,p}(\mathsf W_n,u)^2.
\label{eq:split-square}
\end{equation}
Apply the power-series logarithm normalized by \(\Logz(1)=0\).
Multiplication by the C51 exponent \(2/n\) yields
\[
\frac2n\Logz L_{K,p}
=\frac4n\Logz L_{\Q,p}.
\]
The remaining denominator is \(n/\gcd(n,4)\).  Hence the third row
retains exponent \(4/3\), while the fourth row has exponent one.
Exponentiating in the same origin-normalized branch proves
\eqref{eq:split-repair}.

\begin{remark}[Meaning of ordinary]
The fourth-row result is ordinary in rank and local multiplicity:
one rank-\(255\) rational factor occurs with exponent one.  The packet
contains Tate-normalized pieces, so its polynomial is naturally in
\(\Q[T]\) and need not be integral.  The integral statement in
Theorem~\ref{thm:compatible-main} concerns the untwisted rank-\(10\)
smooth-projective summand.
\end{remark}

Theorem~\ref{thm:packets-repair} now follows.  For \(n\ge5\), the same
denominator arithmetic is only conditional on smoothness and an analogous
source-packet extraction.  The all-\(n\) theorem proved in this paper
remains the equation descent.
